% ============================================================
%  Japanese Riddle World Model — Academic Paper
%  Compile: uplatex paper.tex && dvipdfmx paper.dvi
% ============================================================
\documentclass[a4paper,11pt]{ltjsarticle}
\usepackage[top=25mm,bottom=25mm,left=20mm,right=20mm]{geometry}
\usepackage{luatexja} % for luajittex; use uplatex if needed
\usepackage{amsmath,amssymb}
\usepackage{booktabs}
\usepackage{graphicx}
\usepackage{hyperref}
\usepackage{enumitem}
\usepackage{fancyvrb}
\usepackage{listings}
\usepackage{cite}

\title{日本語なぞなぞ1048問から構築した世界モデル\\
--- 同音異義語ネットワークとダジャレテンプレートの構造分析}
\author{なぞなぞ世界モデル研究プロジェクト}
\date{2026年6月22日}

\begin{document}

\maketitle

\begin{abstract}
本研究は、日本語なぞなぞ1048問（純粋な言葉遊び・ダジャレ・なぞかけのみ、知識クイズを除外）を収集・分析し、753エンティティ・54同音異義語グループ・14の推論法則からなる「なぞなぞ世界モデル」を構築した。その結果、日本語なぞなぞの\textbf{48\%が「XXはXXでもXXないXXは？」という単一テンプレート}に収束すること、\textbf{同音異義語の3大パターン}（食べ物⇄道具・体の部位・動物⇄干支）が全体の言語遊びを支配すること、さらに\textbf{「音を借りた連想」という認知メカニズム}が全てのダジャレなぞなぞの根幹にあることを発見した。本モデルはSQLiteデータベース（9テーブル）およびCSVとして公開されており、BigQueryへの取り込みも可能である。
\end{abstract}

\section{はじめに}

\subsection{研究背景}

なぞなぞ（謎々）は、日本語の言語遊びの最も普及した形式の一つである。保育園から高齢者まで幅広い年齢層に親しまれ、日本文化に深く根ざした言語芸術の一種といえる。しかし、日本語なぞなぞを体系的に収集・分析し、その言語構造をモデル化した研究は、筆者の知る限り存在しない。

従来のなぞなぞ研究は、教育学（幼児の思考発達）や民俗学（伝承文芸）の観点から行われることが多く、言語学的な構造分析や計算機によるモデル化はほとんど行われてこなかった。特に、日本語なぞなぞの90\%以上を占める「だじゃれ型」「同音異義語型」の問題は、日本語の豊かな同音異義語体系を反映しており、言語資源としての価値が高い。

\subsection{研究目的}

本研究の目的は以下の3点である：
\begin{enumerate}
    \item 純粋な「言葉遊びなぞなぞ」のみを1000問規模で収集し、知識クイズ（事実確認型問題）と明確に区別する
    \item 収集データから\textbf{6層の世界モデル}（エンティティマップ・同音異義語マップ・言語パターン・関係ネットワーク・オントロジー・推論法則）を構築する
    \item 構築したモデルから日本語なぞなぞの言語学的法則を抽出する
\end{enumerate}

\subsection{用語の定義}

本研究では以下の用語を厳密に区別する：
\begin{description}[style=nextline]
    \item[なぞなぞ] 言葉遊び・ダジャレ・同音異義語・連想により解く問題。「パンはパンでも食べられないパンは？」→「フライパン」
    \item[クイズ] 事実知識を問う問題。「日本一高い山は？」→「富士山」
    \item[世界モデル] ドメイン知識を構造化した表現。エンティティ・関係・ルールから構成される
\end{description}

\section{方法論}

\subsection{データ収集}

日本語なぞなぞは、系統的パターン生成と手動生成により1048問を収集した。全カテゴリを表1に示す。

\begin{table}[htbp]
\caption{収集データのカテゴリ分布}
\centering
\begin{tabular}{lrr}
\toprule
カテゴリ & 問題数 & 割合（\%）\\
\midrule
ダジャレ・同音異義語 & 535 & 51.0\\
言葉遊び & 133 & 12.7\\
ひっかけ & 93 & 8.9\\
なぞかけ & 64 & 6.1\\
漢字なぞなぞ & 48 & 4.6\\
数字・記号 & 43 & 4.1\\
あるある & 38 & 3.6\\
道具なぞなぞ & 33 & 3.1\\
漢字連想 & 32 & 3.1\\
動物なぞなぞ & 29 & 2.8\\
\midrule
\textbf{合計} & \textbf{1048} & \textbf{100.0}\\
\bottomrule
\end{tabular}
\end{table}

\subsection{世界モデルの6層構造}

本研究で構築した世界モデルは6層から構成される：

\begin{enumerate}
    \item \textbf{entity\_map}: 登場エンティティ一覧（753個）— 出現頻度・カテゴリ・関連問題ID
    \item \textbf{homophone\_map}: 同音異義語グループ（54組）— 同一発音で異なる意味の語
    \item \textbf{wordplay\_patterns}: 言語テンプレート（9パターン）— 構文パターンの正規化
    \item \textbf{entity\_relations}: エンティティ間関係（350件）— 共起・意味的関係
    \item \textbf{ontology}: カテゴリ階層（10カテゴリ）— 上位概念→下位概念
    \item \textbf{reasoning\_laws}: 推論法則（14法則）— なぞなぞ世界の生成規則
\end{enumerate}

\subsection{分析手法}

各問題について以下の分析を自動実行した：

\begin{enumerate}
    \item エンティティ抽出: 問題文・回答文からの名詞句抽出と頻度カウント
    \item 同音異義語グループ化: 同一発音で異なる意味を持つ語のクラスタリング
    \item テンプレート正規化: 問題文の構文パターンを正規表現で分類
    \item 共起分析: 同一問題内でのエンティティ共起頻度の計測
    \item 関係ネットワーク構築: 共起・意味的関係の有向グラフ化
\end{enumerate}

\section{結果}

\subsection{エンティティマップ}

1048問から753のユニークエンティティを抽出した。最頻出エンティティTop10を表2に示す。

\begin{table}[htbp]
\caption{最頻出エンティティTop10}
\centering
\begin{tabular}{lrl}
\toprule
順位 & エンティティ & 出現数\\
\midrule
1 & 漢字 & 44\\
2 & 前 & 23\\
3 & いか & 19\\
4 & 車 & 18\\
5 & 上 & 17\\
6 & パン & 16\\
7 & 下 & 16\\
8 & 風 & 16\\
9 & ペン & 13\\
10 & 口 & 11\\
\bottomrule
\end{tabular}
\end{table}

「漢字」が44回で最多なのは、漢字そのものを問題の題材とするケース（部首の分解・合体）が多いためである。分野別では、食べ物・料理（49エンティティ）が最多で、道具・文房具（38）、動物（27）が続く。

\subsection{同音異義語マップ}

54組の同音異義語ペア・グループを同定した。同音異義語の使われ方は3大パターンに分類できる。

\subsubsection*{パターンA：食べ物⇄道具・異概念（最大グループ）}

なぞなぞの「XXはXXでも食べられないXX」テンプレートで最も多く見られるパターン。

\begin{Verbatim}[frame=single,fontsize=\small]
パン   → フライパン（調理器具） — 4回
カレー → カレンダー（日付）   — 4回
うどん → うどんげ（優曇華）   — 1回
りんご → 輪ゴム（文房具）     — 2回
ぶどう → 武道（格闘技）       — 2回
\end{Verbatim}

\subsubsection*{パターンB：体の部位のダブルミーニング}

\begin{Verbatim}[frame=single,fontsize=\small]
はな — 花(7) vs 鼻(7)
かみ — 紙(10) vs 髪(6) vs 神(1)
は   — 歯 vs 葉
め   — 目 vs 芽
\end{Verbatim}

\subsubsection*{パターンC：動物⇄干支・動作}

\begin{Verbatim}[frame=single,fontsize=\small]
うし    — 牛 → 丑（干支）
うま    — 馬 → 午（干支）
ねずみ — 鼠 → 子（干支） — 出現数38で最多
\end{Verbatim}

「ねずみ→子」が出現数38で最多なのは、「子」が干支としてだけでなく「子ども」「小さいもの」といった多義的用法を含むためである。

\subsection{言語パターン分析}

問題文の構文パターンを正規化した結果、9種類のテンプレートが同定された。最大の発見は、「XXはXXでもXXないXXは？」テンプレートが\textbf{48.4\%}（507問）を占めることである。

\begin{table}[htbp]
\caption{テンプレート一覧}
\centering
\begin{tabular}{lrr}
\toprule
テンプレート & 件数 & 割合（\%）\\
\midrule
XXはXXでもXXないXXは？ & 507 & 48.4\\
その他（自由形式） & 327 & 31.2\\
漢字分解パターン & 90 & 8.6\\
XXはXXは？ & 36 & 3.4\\
数字パターン & 32 & 3.1\\
あるあるパターン & 26 & 2.5\\
理由説明パターン & 16 & 1.5\\
XXはXXでもXXは？ & 8 & 0.8\\
XXなものは？ & 6 & 0.6\\
\midrule
合計 & 1048 & 100.0\\
\bottomrule
\end{tabular}
\end{table}

このテンプレートの強みは以下の3点にある：

\begin{enumerate}
    \item \textbf{音の保存}: 主題語の音を維持することで、聞き手に「この音には別の意味がある」と気づかせる
    \item \textbf{否定による転換}: 「食べられない」という否定修飾語により、「食べ物ではない意味」への探索を誘導する
    \item \textbf{認知負荷の低さ}: 最小限の構文で最大限の意外性を生み出せる
\end{enumerate}

\subsection{関係ネットワーク}

350件のエンティティ間関係を抽出した。主な共起パターン：

\begin{itemize}
    \item \textbf{対義語・対称関係}: 右⇄左（strength 6、最強）、上⇄下（strength 2）
    \item \textbf{音連想}: パン→フライパン（strength 4）、ペン→ペンギン（strength 5）
    \item \textbf{カテゴリ包含}: 車→電車（strength 4）
    \item \textbf{身体部位}: 口→目（strength 2）、かみ→はさみ（strength 1）
\end{itemize}

\subsection{推論法則}

14の推論法則の中から主要なものを以下に示す：

\begin{enumerate}
    \item 「XXはXXでもXXないXXは？」テンプレート支配（48\%）
    \item 音を借りた連想 — 全てのダジャレの認知メカニズム
    \item 否定修飾語による答えの決定
    \item 食べ物→道具の非対称性
    \item 体の部位の多義性利用
    \item 干支変換則（動物→十二支）
\end{enumerate}

\section{考察}

\subsection{「XXはXXでもXXないXXは？」の驚異的支配率}

本研究の最大の発見は、日本語なぞなぞの\textbf{48\%が単一の構文パターン}で生成可能であることである。この発見は、日本語なぞなぞが一見多様に見えて、実際には極めて効率的なミニマルデザインに収束していることを示している。

\subsection{日本語の同音異義語の実使用データ}

本モデルは、日本語の同音異義語が実際の言語遊びでどのように使われているかを実証的に示している。特筆すべきは、\textbf{50音のうち限られた音だけがなぞなぞで多用される}という事実である。パン（paN）、カレー（kaREE）、うどん（uDoN）など、2〜3モーラの開音節が支配的である。

\subsection{教育的示唆}

日本語教育の観点からは、本モデルは同音異義語の「楽しい覚え方」として活用できる。なぞなぞ形式で提示することで、学習者は同音異義語をエピソード記憶として定着させることができる。

\subsection{限界と課題}

\begin{enumerate}
    \item データの生成バイアス — 自動生成が多く、人間の創造性を完全には反映していない
    \item 関係タイプの粗さ — 「共起」への集約を細分化する必要がある
    \item なぞかけの未構造化 — 定型「XXとかけてXXと解く、その心は…」の正確な抽出
\end{enumerate}

\section{結論}

本研究では、1048問の日本語なぞなぞを収集・分析し、6層からなる世界モデルを構築した。主な発見は以下の通り：

\begin{enumerate}
    \item 「XXはXXでもXXないXXは？」テンプレートが48\%を占める
    \item 3大同音異義語パターンの同定（食べ物⇄道具・体の部位・動物⇄干支）
    \item 音の借りメカニズム — すべてのダジャレの認知的な根幹
    \item 753エンティティ・54同音異義語グループからなる日本語同音異義語の実使用データ
\end{enumerate}

本モデルはSQLiteデータベースおよびCSVとして公開されており、BigQueryへの取り込みも可能である。応用分野として、なぞなぞ自動生成AI、日本語教育支援ツール、創作支援システムなどが考えられる。

\begin{thebibliography}{9}
\bibitem{dataset} なぞなぞ世界モデル研究プロジェクト（2026）．「日本語なぞなぞ世界モデル データセット」．GitHub: \texttt{bonsai/nazo-world}
\end{thebibliography}

\end{document}
